\chapter{Inequalities, Large Numbers and Entropy}

\begin{center}\itshape
Jacob Bernoulli wrote that he had pondered it for twenty years: draw often enough from an urn, and the\\
proportion of white balls drawn will, with near certainty, be close to the true proportion. The theorem\\
appeared in 1713, eight years after his death. How often is ``often enough''?
\end{center}

An average of many independent random quantities is almost never far from its expectation. This chapter makes
the words ``almost never'' and ``far'' precise, with three tools of increasing strength. Markov's inequality uses only the
mean; Chebyshev's uses the variance and gives Bernoulli's theorem, the \emph{weak law of large numbers}; and
Chernoff's method, which applies Markov's inequality to an exponential, shows that the probability of a deviation
decays \emph{exponentially} with the number of trials, at a rate that we compute exactly for coin flips.

Two applications show that these crude-looking inequalities have teeth. Bernstein used the weak law to prove
Weierstrass's theorem, that every continuous function on $[0,1]$ is a uniform limit of polynomials, with explicit
polynomials. And the Borel--Cantelli lemmas turn exponential bounds into statements about one infinite sequence of
trials, such as the strong law of large numbers for bounded variables, or the length of the longest run of heads.

The last two sections are about information. The \emph{entropy} of a random variable is the number of bits needed
to describe it, in two precise senses: no prefix code can do better on average, and among $n$ independent copies
almost all the probability sits on about $2^{nH}$ ``typical'' sequences. Jensen's inequality, promised in Volume~0,
is the engine.

All random variables in this chapter are discrete, as in Chapter~2.

\section{Markov and Chebyshev}

\begin{think}
In a town, the mean income is $30\,000$ a year and nobody has a negative income. What fraction of the inhabitants
can earn at least $300\,000$? What if you also know that the standard deviation is $10\,000$?
\end{think}

\begin{proposition}\label{prop:markov}
\begin{enumerate}[label=(\alph*)]
\item \emph{Markov's inequality}: if $X\ge0$ and $a>0$, then $P(X\ge a)\le\dfrac{\mathbb E[X]}a$.
\item \emph{Chebyshev's inequality}: if $\mathbb E[X^2]<\infty$ and $a>0$, then $P\bigl(|X-\mathbb EX|\ge a\bigr)\le\dfrac{\operatorname{Var}X}{a^2}$.
\end{enumerate}
\end{proposition}
\begin{proof}
(a) Pointwise, $a\mathbf 1_{X\ge a}\le X$: if $X\ge a$ the left side is $a\le X$, otherwise it is $0\le X$. Take expectations
(Theorem~\ref{thm:linearity}(b)): $aP(X\ge a)\le\mathbb E[X]$.

(b) Apply (a) to $(X-\mathbb EX)^2$ and $a^2$: the events $\{|X-\mathbb EX|\ge a\}$ and $\{(X-\mathbb EX)^2\ge a^2\}$ are the same.
\end{proof}

So at most a tenth of the town earns ten times the mean income, and if the standard deviation is $10\,000$ then
at most $(10\,000/270\,000)^2\approx0.14\%$ does. Both inequalities are sharp (Exercise~\ref{ex:sharp}): the
first is an equality when $X$ takes only the values $0$ and $a$. They are weak in typical situations only because
they use so little information, and the rest of the chapter feeds them more.

\section{The weak law and Bernstein polynomials}

\begin{theorem}[weak law of large numbers]\label{thm:weaklaw}
Let $X_1,X_2,\ldots$ have finite second moments, the same mean $\mu$, variances at most $\sigma^2$, and be
\emph{uncorrelated}: $\operatorname{Cov}(X_i,X_j)=0$ for $i\ne j$ (for instance, pairwise independent). Let $S_n=X_1+\dots+X_n$. Then for every $\varepsilon>0$,
\[
P\Bigl(\Bigl|\frac{S_n}n-\mu\Bigr|\ge\varepsilon\Bigr)\le\frac{\sigma^2}{n\varepsilon^2}\xrightarrow[n\to\infty]{}0.
\]
\end{theorem}
\begin{proof}
By Proposition~\ref{prop:varsum}, $\operatorname{Var}S_n=\sum_i\operatorname{Var}X_i\le n\sigma^2$, so $\operatorname{Var}\frac{S_n}n\le\frac{\sigma^2}n$, while $\mathbb E\frac{S_n}n=\mu$ by
linearity. Apply Chebyshev's inequality.
\end{proof}

This is Bernoulli's theorem: for independent trials with success probability $p$, the frequency of success
$\frac{S_n}n$ is within $\varepsilon$ of $p$ with probability at least $1-\frac{p(1-p)}{n\varepsilon^2}\ge1-\frac1{4n\varepsilon^2}$. We say that $\frac{S_n}n$ converges
to $p$ \emph{in probability}. The bound is honest but pessimistic: to guarantee $99\%$ confidence of being within
$0.01$ of $p$, it asks for $n=250\,000$ trials, while Section~\ref{sec:chernoff} will show that $n=26\,500$ are enough.

Bernstein noticed in 1912 that this theorem, read backwards, approximates functions. If $S_n$ is binomial$(n,x)$, the
frequency $\frac{S_n}n$ is close to $x$, so $f(\frac{S_n}n)$ is close to $f(x)$ for a continuous $f$, and its expectation is a
polynomial in $x$.

\begin{definition}
The \emph{Bernstein polynomials} of a function $f:[0,1]\to\R$ are
\[
\begin{aligned}B_nf(x)&=\sum_{k=0}^nf\Bigl(\frac kn\Bigr)\binom nkx^k(1-x)^{n-k}\eqbr=\mathbb E\Bigl[f\Bigl(\frac{S_n}n\Bigr)\Bigr],\end{aligned}
\]
where $S_n$ is binomial$(n,x)$.
\end{definition}

The two expressions agree by the transfer formula. We need one fact about continuous functions that Volume~0 used
without proof.

\begin{lemma}[Heine]\label{lem:heine}
A continuous function $f:[0,1]\to\R$ is \emph{uniformly continuous}: for every $\varepsilon>0$ there is $\delta>0$ such that
$|x-y|<\delta$ implies $|f(x)-f(y)|<\varepsilon$. In particular $f$ is bounded.
\end{lemma}
\begin{proof}
Suppose not: for some $\varepsilon>0$ and every $n$ there are $x_n,y_n\in[0,1]$ with $|x_n-y_n|<\frac1n$ and $|f(x_n)-f(y_n)|\ge\varepsilon$. By
Bolzano--Weierstrass (Volume~0, Proposition~5.1(c)) some subsequence $x_{n_k}$ converges to some $x\in[0,1]$, and then
$y_{n_k}\to x$ as well. By continuity at $x$, $f(x_{n_k})-f(y_{n_k})\to f(x)-f(x)=0$, a contradiction. For boundedness, take
$\varepsilon=1$ and the corresponding $\delta$: any $x\in[0,1]$ is reached from $0$ in $N=\lceil2/\delta\rceil$ equal steps of length $\frac xN<\delta$, and $f$
changes by less than $1$ at each step, so $|f(x)|\le|f(0)|+N$.
\end{proof}

\begin{theorem}[Weierstrass, Bernstein's proof]\label{thm:weierstrass}
If $f:[0,1]\to\R$ is continuous, then $B_nf\to f$ uniformly on $[0,1]$. In particular, every continuous function on
$[0,1]$ is a uniform limit of polynomials.
\end{theorem}
\begin{proof}
Let $M=\sup|f|$, finite by Lemma~\ref{lem:heine}. Fix $\varepsilon>0$ and take $\delta$ as in the lemma. For $x\in[0,1]$, let $A$ be the event
$\bigl\{|\frac{S_n}n-x|\ge\delta\bigr\}$. On $A^c$, $|f(\frac{S_n}n)-f(x)|<\varepsilon$; on $A$, it is at most $2M$. So
\[
\begin{aligned}
|B_nf(x)-f(x)|&\le\mathbb E\Bigl|f\Bigl(\frac{S_n}n\Bigr)-f(x)\Bigr|\eqbr\le\varepsilon+2M\,P(A)\\
&\le\varepsilon+2M\,\frac{x(1-x)}{n\delta^2}\le\varepsilon+\frac M{2n\delta^2},
\end{aligned}
\]
by Chebyshev, since $\operatorname{Var}\frac{S_n}n=\frac{x(1-x)}n\le\frac1{4n}$. The bound does not depend on $x$, and it is less than $2\varepsilon$ once
$n>\frac M{2\varepsilon\delta^2}$.
\end{proof}

\begin{figure}[ht]
\centering
\fig{ch04-bernstein}
\caption{The Bernstein polynomials $B_nf$ of $f(x)=|x-\frac12|$ (black) for $n=4$, $16$ and $64$. They lie above $f$, since $f$ is
convex (this is Jensen's inequality, Section~\ref{sec:entropy}), and the worst error, at the kink, is of order $\frac1{\sqrt n}$: here
$0.19$, $0.10$ and $0.05$.}
\label{fig:bernstein}
\end{figure}

\begin{beauty}
Weierstrass's 1885 proof used integrals against Gaussian kernels; Bernstein's proof uses a coin. The polynomial
$B_nf$ is the answer to a betting question: toss a coin of bias $x$ a total of $n$ times, and receive $f$ of the
observed frequency. Its expected payoff is a polynomial in $x$, and the law of large numbers says that it is close
to $f(x)$. Pierre B\'ezier rediscovered the same polynomials in the 1960s to design car bodies at Renault, and every
smooth curve on a computer screen is drawn with them today.
\end{beauty}

\section{Chernoff bounds}\label{sec:chernoff}

\begin{think}
Toss a fair coin $1000$ times. Chebyshev says that the probability of at least $600$ heads is at most
$\frac{250}{100^2}=0.025$. Is the truth closer to $10^{-2}$, $10^{-5}$ or $10^{-10}$?
\end{think}

Chebyshev applies Markov to $(X-\mu)^2$. Chernoff's idea is to apply it to $e^{\lambda X}$ and then choose the best $\lambda$: for
every $\lambda>0$,
\[
P(X\ge a)=P\bigl(e^{\lambda X}\ge e^{\lambda a}\bigr)\le e^{-\lambda a}\,\mathbb E\bigl[e^{\lambda X}\bigr],
\]
and for a sum of independent variables, the \emph{moment generating function} $\mathbb E[e^{\lambda S}]$ is a product.

\begin{theorem}[Chernoff bounds]\label{thm:chernoff}
Let $X_1,\dots,X_n$ be independent, $X_i$ Bernoulli$(p_i)$, $S=\sum X_i$ and $\mu=\mathbb E[S]=\sum p_i>0$. For $\delta>0$ and $0<\delta'<1$,
\[
\begin{gathered}
P\bigl(S\ge(1+\delta)\mu\bigr)\le\Bigl(\frac{e^\delta}{(1+\delta)^{1+\delta}}\Bigr)^{\mu},\qq P\bigl(S\le(1-\delta')\mu\bigr)\le\Bigl(\frac{e^{-\delta'}}{(1-\delta')^{1-\delta'}}\Bigr)^{\mu}.
\end{gathered}
\]
Consequently, for $0<\delta\le1$, $P(S\ge(1+\delta)\mu)\le e^{-\delta^2\mu/3}$, $P(S\le(1-\delta)\mu)\le e^{-\delta^2\mu/2}$, and
$P(|S-\mu|\ge\delta\mu)\le2e^{-\delta^2\mu/3}$.
\end{theorem}
\begin{proof}
For real $\lambda$, the variables $e^{\lambda X_i}$ are independent (Proposition~\ref{prop:indepfun}), so by Theorem~\ref{thm:linearity}(c) and induction,
\[
\begin{aligned}\mathbb E[e^{\lambda S}]&=\prod_{i=1}^n\bigl(1+p_i(e^\lambda-1)\bigr)\eqbr\le\prod_{i=1}^ne^{p_i(e^\lambda-1)}=e^{\mu(e^\lambda-1)},\end{aligned}
\]
using $1+u\le e^u$. Upper tail: with $\lambda=\ln(1+\delta)>0$, Markov gives
$P(S\ge(1+\delta)\mu)\le e^{-\lambda(1+\delta)\mu}e^{\mu\delta}=\bigl(e^\delta(1+\delta)^{-(1+\delta)}\bigr)^\mu$. Lower tail: with $\lambda=\ln(1-\delta')<0$,
\[
\begin{aligned}
P\bigl(S\le(1-\delta')\mu\bigr)&=P\bigl(e^{\lambda S}\ge e^{\lambda(1-\delta')\mu}\bigr)\\
&\le e^{-\lambda(1-\delta')\mu}\,e^{\mu(e^\lambda-1)}\eqbr=\bigl(e^{-\delta'}(1-\delta')^{-(1-\delta')}\bigr)^\mu.
\end{aligned}
\]
For the simplified forms, take logarithms. The function $f(\delta)=\delta-(1+\delta)\ln(1+\delta)+\frac{\delta^2}3$ has $f(0)=0$,
$f'(\delta)=-\ln(1+\delta)+\frac{2\delta}3$, $f'(0)=0$ and $f''(\delta)=\frac23-\frac1{1+\delta}$, negative on $[0,\frac12[$ and positive after. So $f'$
decreases and then increases on $[0,1]$, and $f'\le\max(f'(0),f'(1))=\max(0,\frac23-\ln2)=0$; hence $f\le0$ on $[0,1]$, which is
the first bound. Likewise $g(\delta)=-\delta-(1-\delta)\ln(1-\delta)+\frac{\delta^2}2$ has $g(0)=0$ and $g'(\delta)=\ln(1-\delta)+\delta\le0$, so $g\le0$; this
gives the second bound for $\delta<1$, and for $\delta=1$ the bound $P(S\le0)=\prod(1-p_i)\le e^{-\mu}\le e^{-\mu/2}$ is direct. The third bound is
the union of the first two, since $e^{-\delta^2\mu/2}\le e^{-\delta^2\mu/3}$.
\end{proof}

For $1000$ fair coins, $\mu=500$ and $\delta=\frac15$ give $P(S\ge600)\le e^{-500/75}\approx1.3\times10^{-3}$, already twenty times better than
Chebyshev. The simplified bounds are convenient but lossy; for identically distributed trials the method gives the
exact exponential rate, as we now show. Write, for $0<p<1$ and $0\le a\le1$,
\[
I_p(a)=a\ln\frac ap+(1-a)\ln\frac{1-a}{1-p}\qquad(0\ln0=0).
\]

\begin{theorem}[large deviations for coin flips]\label{thm:cramer}
Let $X_1,X_2,\ldots$ be independent Bernoulli$(p)$, $S_n=X_1+\dots+X_n$ and $p<a<1$. Then for every $n\ge1$
\[
\frac1{n+1}e^{-nI_p(\lceil na\rceil/n)}\le P(S_n\ge na)\le e^{-nI_p(a)},
\]
and consequently $\frac1n\ln P(S_n\ge na)\to-I_p(a)$. By symmetry ($S_n\mapsto n-S_n$, $p\mapsto 1-p$),
$P(S_n\le na)\le e^{-nI_p(a)}$ for $0<a<p$.
\end{theorem}
\begin{proof}
\emph{Upper bound.} Here $\mathbb E[e^{\lambda S_n}]=(1-p+pe^\lambda)^n$, so for $\lambda>0$, $P(S_n\ge na)\le\exp\bigl(-n(\lambda a-\ln(1-p+pe^\lambda))\bigr)$.
Write $q=1-p$, $b=1-a$, and choose $e^\lambda=\frac{aq}{pb}$, which is $>1$ since $a>p$ and $b<q$. Then $1-p+pe^\lambda=q+\frac{aq}b=\frac qb$ and
\[
\begin{aligned}\lambda a-\ln\frac qb&=a\ln\frac ap+a\ln\frac qb-\ln\frac qb\eqbr=a\ln\frac ap+b\ln\frac bq=I_p(a).\end{aligned}
\]
\emph{Lower bound.} Let $k=\lceil na\rceil$ and $\alpha=\frac kn$, so $a\le\alpha$. If $k=n$, $P(S_n\ge na)=p^n=e^{-nI_p(1)}$ and we are done; so let
$k<n$, i.e.\ $0<\alpha<1$. The binomial$(n,\alpha)$ probabilities $t_j=\binom nj\alpha^j(1-\alpha)^{n-j}$, $0\le j\le n$, add up to $1$, and the largest is $t_k$:
indeed $\frac{t_{j+1}}{t_j}=\frac{(n-j)\alpha}{(j+1)(1-\alpha)}$, which is $\ge1$ iff $j+1\le(n+1)\alpha=k+\alpha$, i.e.\ iff $j<k$. So $t_k\ge\frac1{n+1}$. Now
\[
\begin{aligned}P(S_n\ge na)&\ge P(S_n=k)\eqbr=t_k\Bigl(\frac p\alpha\Bigr)^k\Bigl(\frac q{1-\alpha}\Bigr)^{n-k}=t_k\,e^{-nI_p(\alpha)}.\end{aligned}
\]
\emph{Limit.} $\frac1n\ln(n+1)\to0$, and $\frac{\lceil na\rceil}n\to a$, so $I_p(\lceil na\rceil/n)\to I_p(a)$ by continuity of $I_p$ on $]0,1[$.
\end{proof}

So the probability of a deviation is neither ``small'' nor ``polynomially small'': it is $e^{-nI_p(a)+o(n)}$, with an explicit
exponent. For $1000$ fair coins and $a=0.6$, $I_{1/2}(0.6)\approx0.0201$ and the upper bound is $1.8\times10^{-9}$; the exact value is
$1.36\times10^{-10}$, and the factor of $13$ between them is the polynomial correction that the theorem does not see. The
truth was $10^{-10}$: Chebyshev was off by a factor of a hundred million.

The quantity $I_p(a)$ is a \emph{relative entropy}, measured in nats rather than bits (Section~\ref{sec:entropy}). It is not
symmetric around $p$, and it is at least $2(a-p)^2$ (Exercise~\ref{ex:pinsker}), which gives the cruder but universal bound
of the next theorem.

\begin{lemma}[Hoeffding]\label{lem:hoeffding}
Let $X$ take values in $[0,1]$, with mean $p$. For every real $\lambda$, $\mathbb E\bigl[e^{\lambda(X-p)}\bigr]\le e^{\lambda^2/8}$.
\end{lemma}
\begin{proof}
The function $x\mapsto e^{\lambda x}$ is convex, so on $[0,1]$ it lies below its chord: $e^{\lambda x}\le(1-x)+xe^\lambda$. Taking expectations,
$\mathbb E[e^{\lambda X}]\le1-p+pe^\lambda$. Let $L(\lambda)=\ln(1-p+pe^\lambda)-\lambda p$; we show $L(\lambda)\le\frac{\lambda^2}8$. Now $L(0)=0$, and with
$u(\lambda)=\frac{pe^\lambda}{1-p+pe^\lambda}\in[0,1]$ we get $L'=u-p$, so $L'(0)=0$, and $L''=u(1-u)\le\frac14$. Taylor's formula with integral remainder (Volume~0,
Chapter~6) gives $L(\lambda)=\int_0^\lambda(\lambda-t)L''(t)\dd t\le\frac14\cdot\frac{\lambda^2}2$ (for $\lambda<0$ as well, since $(\lambda-t)\dd t$ is then also of
constant sign along the path and the integral equals $\int_\lambda^0(t-\lambda)L''(t)\dd t$).
\end{proof}

\begin{theorem}[Hoeffding's inequality]\label{thm:hoeffding}
Let $X_1,\dots,X_n$ be independent with values in $[a_i,b_i]$, $S=\sum X_i$, and $t>0$. Then
\[
P(S-\mathbb ES\ge t)\le\exp\Bigl(-\frac{2t^2}{\sum_i(b_i-a_i)^2}\Bigr),
\]
and the same bound holds for $P(S-\mathbb ES\le-t)$. For $X_i$ in $[0,1]$, the bound is $e^{-2t^2/n}$.
\end{theorem}
\begin{proof}
Let $c_i=b_i-a_i$ (we may assume $c_i>0$, a constant $X_i$ contributing nothing), and $Y_i=\frac{X_i-a_i}{c_i}\in[0,1]$. By the lemma applied to
$Y_i$ with $\lambda c_i$ in place of $\lambda$, $\mathbb E[e^{\lambda(X_i-\mathbb EX_i)}]=\mathbb E[e^{\lambda c_i(Y_i-\mathbb EY_i)}]\le e^{\lambda^2c_i^2/8}$. By independence,
$\mathbb E[e^{\lambda(S-\mathbb ES)}]\le e^{\lambda^2C/8}$ with $C=\sum c_i^2$, and by Markov, for $\lambda>0$,
\[
P(S-\mathbb ES\ge t)\le e^{-\lambda t+\lambda^2C/8}.
\]
The exponent is smallest for $\lambda=\frac{4t}C$, where it equals $-\frac{2t^2}C$. For the lower tail, apply this to $-X_i\in[-b_i,-a_i]$.
\end{proof}

With $t=\varepsilon n$ and $X_i\in[0,1]$: $P(|\frac{S_n}n-\mu|\ge\varepsilon)\le2e^{-2n\varepsilon^2}$. For $99\%$ confidence of being within $0.01$, $26\,500$ trials suffice
(since $2e^{-2n\cdot10^{-4}}\le0.01$ iff $n\ge5000\ln200\approx26\,492$), against $250\,000$ for Chebyshev. Behind every
``margin of error'' and every Monte Carlo error bar that comes with a guarantee, there is an inequality of this kind. Chapter~10 extends it to martingales
(Azuma's inequality), and Volume~10 studies large deviations in general (Cram\'er's theorem).

\begin{figure}[ht]
\centering
\fig{ch04-chernoff}
\caption{The probability that $n$ trials with success probability $p=0.1$ give a success frequency of at least $0.2$ (black dots,
computed exactly), against Chebyshev's bound $\frac{9}{n}$ (orange), Hoeffding's bound $e^{-2n(0.1)^2}$ (green) and the optimal
Chernoff bound $e^{-nI_{0.1}(0.2)}$ (blue), on a logarithmic scale. Chebyshev decays like a power of $n$; the others are straight
lines, and the blue one has the true slope $-I_{0.1}(0.2)\approx-0.044$, more than twice Hoeffding's $-0.02$.}
\label{fig:chernoff}
\end{figure}

\section{Borel--Cantelli}\label{sec:bc}

The weak law speaks about one $n$ at a time. To say something about a whole infinite sequence of trials (whose
existence is admitted, as in Chapter~1) we need events such as ``$A_n$ happens for infinitely many $n$''.

\begin{definition}
For events $A_1,A_2,\ldots$, the event \emph{$A_n$ infinitely often} is
\[
\{A_n\text{ i.o.}\}=\bigcap_{N\ge1}\bigcup_{n\ge N}A_n,
\]
the set of $\omega$ that belong to infinitely many $A_n$. Its complement, $\bigcup_N\bigcap_{n\ge N}A_n^c$, is ``eventually $A_n^c$''. A property holds
\emph{almost surely} (a.s.) if the event where it holds has probability $1$.
\end{definition}

\begin{think}
A monkey hits the keys of a typewriter at random, forever. Will it type the word \texttt{banana}? Infinitely
often? What about the complete works of Shakespeare?
\end{think}

\begin{theorem}[Borel--Cantelli lemmas]\label{thm:bc}
Let $A_1,A_2,\ldots$ be events.
\begin{enumerate}[label=(\alph*)]
\item If $\sum_nP(A_n)<\infty$, then $P(A_n\text{ i.o.})=0$.
\item If the $A_n$ are independent and $\sum_nP(A_n)=\infty$, then $P(A_n\text{ i.o.})=1$.
\end{enumerate}
\end{theorem}
\begin{proof}
(a) For every $N$, $\{A_n\text{ i.o.}\}\subseteq\bigcup_{n\ge N}A_n$, so by the union bound $P(A_n\text{ i.o.})\le\sum_{n\ge N}P(A_n)$, the tail of a
convergent series, which tends to $0$.

(b) We show $P\bigl(\bigcap_{n\ge N}A_n^c\bigr)=0$ for every $N$; the complement of $\{A_n\text{ i.o.}\}$ is the countable union of these events. The
$A_n^c$ are independent (Proposition~\ref{prop:indepcomp}), so for $M\ge N$, using $1-u\le e^{-u}$,
\[
\begin{aligned}P\Bigl(\bigcap_{n=N}^MA_n^c\Bigr)&=\prod_{n=N}^M\bigl(1-P(A_n)\bigr)\eqbr\le\exp\Bigl(-\sum_{n=N}^MP(A_n)\Bigr)\xrightarrow[M\to\infty]{}0.\end{aligned}
\]
These events decrease in $M$ to $\bigcap_{n\ge N}A_n^c$, which therefore has probability $0$ by continuity (Proposition~\ref{prop:rules}(e)).
\end{proof}

The monkey types \texttt{banana} infinitely often: cut its output into consecutive blocks of six keys; the events
``block $j$ spells \texttt{banana}'' are independent with the same positive probability, so their probabilities sum to $\infty$.
The same holds for Shakespeare, although the waiting time is beyond astronomical. Without independence, (b) fails:
take all $A_n$ equal to one event of probability $\frac12$.

The first lemma, fed with Hoeffding's inequality, gives a law of large numbers for a single infinite sequence.

\begin{corollary}[strong law for bounded variables]\label{cor:slln}
Let $X_1,X_2,\ldots$ be independent discrete random variables with values in $[a,b]$ and the same mean $\mu$. Then
$\frac{S_n}n\to\mu$ almost surely.
\end{corollary}
\begin{proof}
Fix $\varepsilon>0$. By Hoeffding, $P(|\frac{S_n}n-\mu|\ge\varepsilon)\le2e^{-2n\varepsilon^2/(b-a)^2}$ (for $b>a$; if $a=b$ there is nothing to prove), the
general term of a convergent geometric series. By Borel--Cantelli (a), the event $E_\varepsilon$ ``eventually $|\frac{S_n}n-\mu|<\varepsilon$'' has
probability $1$. Then $E=\bigcap_{k\ge1}E_{1/k}$ has probability $1$ (its complement is a countable union of null events), and on $E$,
$\frac{S_n}n\to\mu$.
\end{proof}

Chapter~6 proves the strong law under a fourth-moment assumption by the same route, with Chebyshev applied to
the fourth power; the strong law for all integrable i.i.d.\ variables (Kolmogorov, Etemadi) needs a truncation
argument and comes later in the series. With the binary digits of a uniform random number, which Chapter~5 shows to
be independent fair bits, Corollary~\ref{cor:slln} becomes Borel's theorem of 1909: almost every number in $[0,1]$
has, in the limit, as many $0$s as $1$s in its binary expansion.

The two lemmas together locate the longest runs of heads, a test that human-made ``random'' sequences
regularly fail (a numerical calculation has an experiment).

\begin{proposition}\label{prop:runs}
Toss a fair coin infinitely often, and let $\ell_n$ be the length of the run of heads starting at toss $n$ ($\ell_n=0$ if
toss $n$ is tails). Then, almost surely, $\displaystyle\limsup_{n\to\infty}\frac{\ell_n}{\log_2n}=1$.
\end{proposition}
\begin{proof}
For $k\ge0$, $\ell_n\ge k$ iff tosses $n,\dots,n+k-1$ are heads, so $P(\ell_n\ge k)=2^{-k}$.

\emph{At most $1$.} Fix $\varepsilon>0$. Then $P(\ell_n\ge(1+\varepsilon)\log_2n)\le2^{-(1+\varepsilon)\log_2n}=n^{-1-\varepsilon}$, and $\sum_nn^{-1-\varepsilon}<\infty$ (Volume~0,
Chapter~5). By Borel--Cantelli (a), almost surely $\ell_n<(1+\varepsilon)\log_2n$ for all large $n$, hence $\limsup\frac{\ell_n}{\log_2n}\le1+\varepsilon$.
Intersecting over $\varepsilon=\frac1k$, the $\limsup$ is at most $1$ almost surely.

\emph{At least $1$.} We need disjoint windows, to have independence. Let $r_j=\lfloor\log_2j\rfloor$ for $j\ge1$, let $n_1=1$ and
$n_{j+1}=n_j+r_j+1$, and let $B_j$ be the event that tosses $n_j,\dots,n_j+r_j-1$ are all heads. The windows are disjoint, so for
distinct $j_1,\dots,j_m$ the event $B_{j_1}\cap\dots\cap B_{j_m}$ prescribes $r_{j_1}+\dots+r_{j_m}$ distinct tosses and has probability
$2^{-r_{j_1}}\cdots2^{-r_{j_m}}$: the $B_j$ are independent. Since $P(B_j)=2^{-r_j}\ge\frac1j$, Borel--Cantelli (b) says that almost surely
$\ell_{n_j}\ge r_j$ for infinitely many $j$. Finally $n_j\le1+(j-1)(\log_2j+1)\le j(\log_2j+2)$, so $\log_2n_j\le\log_2j+\log_2(\log_2j+2)$, while $r_j\ge\log_2j-1$, and
$\frac{r_j}{\log_2n_j}\to1$. Hence $\limsup\frac{\ell_n}{\log_2n}\ge1$ almost surely.
\end{proof}

So runs of length close to $\log_2n$ keep appearing (Exercise~\ref{ex:longestrun} shows that the longest run among the first $n$ tosses is
$\sim\log_2n$). In $200$ tosses the longest run of heads has mean about $7$, and with probability $97\%$ some run of $6$ identical
outcomes appears. When people are asked to write down $200$ ``random'' tosses, they rarely write a run longer than $5$.

\section{Jensen and entropy}\label{sec:entropy}

Volume~0 proved Jensen's inequality for finitely many points. Here is the form promised there, for random variables.

\begin{theorem}[Jensen]\label{thm:jensenprob}
Let $f$ be twice differentiable and convex on an interval $I$, and $X$ an integrable random variable with values in $I$
such that $f(X)$ is integrable. Then $\mathbb E[X]\in I$ and
\[
f\bigl(\mathbb E[X]\bigr)\le\mathbb E\bigl[f(X)\bigr].
\]
If $f''>0$ on $I$, equality holds only if $X$ is almost surely constant.
\end{theorem}
\begin{proof}
Let $m=\mathbb E[X]$. If $c=\inf I$ is finite, then $X\ge c$ gives $m\ge c$, with equality only if $X=c$ almost surely,
in which case $c$ is a value of $X$ and lies in $I$; the same holds for $\sup I$. So $m\in I$. By Volume~0, Proposition~4.1(b), $f(x)\ge f(m)+f'(m)(x-m)$ for every
$x\in I$. Put $x=X$ and take expectations: $\mathbb E[f(X)]\ge f(m)+f'(m)(\mathbb E[X]-m)=f(m)$. If $f''>0$, the tangent inequality is strict for
$x\ne m$, so $\mathbb E[f(X)]-f(m)$ is the expectation of a non-negative variable that is positive on $\{X\ne m\}$; it vanishes only if
$P(X\ne m)=0$.
\end{proof}

For instance $(\mathbb EX)^2\le\mathbb E[X^2]$, $e^{\mathbb E X}\le\mathbb E[e^X]$, and $\mathbb E[\ln X]\le\ln\mathbb E[X]$ for $X>0$: this last one, with the concave $\ln$, is the
engine of information theory.

\begin{think}
I think of an integer between $1$ and $1000$, and you may ask yes-or-no questions. How many questions do you need?
Now suppose that I choose $1$ half of the time, $2$ a quarter of the time, and so on. How many questions do you
need on average?
\end{think}

\begin{definition}
Let $X$ take values in a countable set $A$, with $p(x)=P(X=x)$. Its \emph{entropy} is
\[
\begin{aligned}H(X)&=\sum_{x\in A}p(x)\log_2\frac1{p(x)}\eqbr=\mathbb E\Bigl[\log_2\frac1{p(X)}\Bigr]\in[0,+\infty],\end{aligned}
\]
with the convention $0\log_2\frac10=0$. For two laws $p$, $q$ on $A$, the \emph{relative entropy} (or Kullback--Leibler divergence)
is $D(p\,\|\,q)=\sum_{x:\,p(x)>0}p(x)\log_2\frac{p(x)}{q(x)}$, which may be $+\infty$.
\end{definition}

The quantity $\log_2\frac1{p(x)}$ measures the surprise of seeing $x$: one bit for a fair coin's outcome, ten bits for
an outcome of probability $2^{-10}$. Entropy is the average surprise. A fair coin has $H=1$; Bernoulli$(p)$ has
$h(p)=p\log_2\frac1p+(1-p)\log_2\frac1{1-p}$, for instance $h(0.1)\approx0.469$; the uniform law on $m$ points has $\log_2m$. The exponent $I_p(a)$ of
Section~\ref{sec:chernoff} is $(\ln2)\,D(\mathrm{B}(a)\,\|\,\mathrm{B}(p))$, where $\mathrm{B}(p)$ denotes the Bernoulli$(p)$ law.

\begin{proposition}[Gibbs' inequality]\label{prop:gibbs}
$D(p\,\|\,q)\ge0$, with equality if and only if $p=q$. Consequently, if $X$ takes at most $m$ values, $H(X)\le\log_2m$, with
equality if and only if $X$ is uniform on $m$ values.
\end{proposition}
\begin{proof}
Let $S=\{x:p(x)>0\}$. Gibbs' inequality is Jensen's inequality $\mathbb E[\ln Y]\le\ln\mathbb E[Y]$ for $Y=\frac{q(X)}{p(X)}$; we run its proof directly, with the
tangent to the concave $\ln$ taken at $1$, which also handles infinite sums and the equality case. By Volume~0,
Proposition~4.1(b), $\ln y\le y-1$ for $y>0$, with equality only at $y=1$. So for $x\in S$ with $q(x)>0$,
\[
p(x)\ln\frac{q(x)}{p(x)}\le q(x)-p(x).
\]
In particular the negative terms of the series defining $D(p\,\|\,q)$ add up to at most $\frac1{\ln2}\sum_xq(x)$, so $D(p\,\|\,q)$ is well defined in
$]-\infty,+\infty]$, and we may assume it finite, so that $q>0$ on $S$. Summing over $S$,
$-(\ln2)\,D(p\,\|\,q)\le\sum_Sq-\sum_Sp=\sum_Sq-1\le0$. Equality forces $\frac{q(x)}{p(x)}=1$ for all $x\in S$, so $q=p$ on $S$, hence everywhere since
both sum to $1$. For the entropy, let $u$ be uniform on a set of $m$ points containing the values of
$X$: $D(p\,\|\,u)=\sum_xp(x)\log_2(m\,p(x))=\log_2m-H(X)\ge0$, with equality iff $p=u$.
\end{proof}

\begin{proposition}\label{prop:entropyadd}
If $X$ and $Y$ are independent with finite entropy, $H(X,Y)=H(X)+H(Y)$. If $X_1,\dots,X_n$ are independent with the law
of $X$, then $H(X_1,\dots,X_n)=nH(X)$.
\end{proposition}
\begin{proof}
$P(X=x,Y=y)=p(x)q(y)$, so $\log_2\frac1{p(x)q(y)}=\log_2\frac1{p(x)}+\log_2\frac1{q(y)}$; multiply by $p(x)q(y)$ and sum over $(x,y)$, a sum of
non-negative terms that may be summed in either order. For $n$ variables, $(X_1,\dots,X_{n-1})$ is independent of $X_n$
(Proposition~\ref{prop:indepfun}); induct.
\end{proof}

\paragraph{Codes.} A \emph{binary code} for a finite or countable alphabet $A$ assigns to each $x\in A$ a word $c(x)\in\{0,1\}^*$ of
length $\ell(x)\ge1$. It is a \emph{prefix code} if no codeword is a prefix of another; then a concatenation of codewords can be
decoded on the fly, left to right, without separators. Morse code is not a prefix code (it needs pauses), while the
codes that compress files are.

\begin{proposition}[Kraft's inequality]\label{prop:kraft}
A prefix code satisfies $\sum_{x\in A}2^{-\ell(x)}\le1$. Conversely, if positive integers $\ell(x)$ satisfy this inequality, there is a
prefix code with these lengths.
\end{proposition}
\begin{proof}
To a word $w=w_1\cdots w_\ell$ associate the \emph{dyadic interval}
\[
I_w=[0.w_1\cdots w_\ell,\;0.w_1\cdots w_\ell+2^{-\ell}[
\]
(binary expansions), of length $2^{-\ell}$: it is the set of $t\in[0,1[$ whose binary expansion (the one not ending with infinitely many $1$s) begins with $w$.
So $I_{w'}\subseteq I_w$ if $w$ is a prefix of $w'$, and $I_w\cap I_{w'}=\varnothing$ if neither word is a prefix of the other, since the
expansions of their points differ at the first place where $w$ and $w'$ differ. For a prefix code the intervals $I_{c(x)}$ are
therefore pairwise disjoint subsets of $[0,1[$, and their total length is at most $1$ (for every finite subset of $A$, then in
the limit).

Conversely, list the letters as $x_1,x_2,\ldots$ with non-decreasing lengths $\ell_i=\ell(x_i)$ (possible even for an infinite alphabet, since the
inequality allows at most $2^\ell$ letters of each length $\ell$), and let $s_i=\sum_{j<i}2^{-\ell_j}$, which is at most $1-2^{-\ell_i}$ by hypothesis. Each $2^{-\ell_j}$ with $j<i$ is a multiple of $2^{-\ell_i}$, so $s_i=k_i2^{-\ell_i}$ for an integer
$0\le k_i<2^{\ell_i}$. Let $c_i$ be the binary expansion of $k_i$ on $\ell_i$ digits; then $I_{c_i}=[s_i,s_{i+1}[$. These intervals are pairwise disjoint,
so no $c_i$ is a prefix of another.
\end{proof}

\begin{theorem}[Shannon's source coding theorem]\label{thm:sourcecoding}
Let $X$ take values in a finite alphabet $A$.
\begin{enumerate}[label=(\alph*)]
\item Every prefix code for $A$ has expected length $\mathbb E[\ell(X)]\ge H(X)$, with equality iff $p(x)=2^{-\ell(x)}$ for all $x$.
\item There is a prefix code with $\mathbb E[\ell(X)]<H(X)+1$, unless $X$ is almost surely constant (then $H(X)=0$ and one bit is
needed).
\item Let $L_n$ be the infimum of the expected lengths of prefix codes for $A^n$, applied to $(X_1,\dots,X_n)$ with $X_i$ independent
copies of $X$. Then $nH(X)\le L_n\le nH(X)+1$, so the optimal number of bits per symbol $\frac{L_n}n$ tends to $H(X)$.
\end{enumerate}
\end{theorem}
\begin{proof}
(a) We may drop the letters with $p(x)=0$, which only decreases $\sum2^{-\ell}$. Let $K=\sum_x2^{-\ell(x)}\le1$ (Kraft) and $q(x)=2^{-\ell(x)}/K$, a law on $A$.
Then
\[
\begin{aligned}\mathbb E[\ell(X)]-H(X)&=\sum_xp(x)\log_2\frac{p(x)}{2^{-\ell(x)}}\eqbr=D(p\,\|\,q)-\log_2K\ge0,\end{aligned}
\]
by Gibbs' inequality and $K\le1$. Equality forces $K=1$ and $p=q$.

(b) If $X$ is constant, one bit is needed and suffices. Otherwise drop the letters of probability $0$, so that $0<p(x)<1$ for every
remaining letter, and let $\ell(x)=\lceil\log_2\frac1{p(x)}\rceil\ge1$. Then $2^{-\ell(x)}\le p(x)$, so $\sum2^{-\ell(x)}\le1$ and Kraft gives a prefix code
with these lengths, and $\ell(x)<\log_2\frac1{p(x)}+1$; multiply by $p(x)$ and sum.

(c) Apply (a) and (b) to the variable $(X_1,\dots,X_n)$ with alphabet $A^n$, whose entropy is $nH(X)$ by Proposition~\ref{prop:entropyadd}.
\end{proof}

In the second question of the last think box, the law $p(k)=2^{-k}$ has entropy $\sum_kk2^{-k}=2$ bits, and the questions
``is it $1$?'', ``is it $2$?'', \ldots\ achieve it: they form the prefix code $1,01,001,\ldots$ with $\ell(k)=k=\log_2\frac1{p(k)}$. For $1000$ equally likely
numbers, ten questions suffice and $\log_21000\approx9.97$ shows that no strategy does better on average than $9.97$.

\begin{remark}
In 1952 David Huffman, then a student of Robert Fano, found an algorithm that produces an \emph{optimal} prefix code:
repeatedly merge the two least likely symbols into one. Its optimality proof is a short exchange argument (Cover and
Thomas, Section~5.8), and a numerical calculation implements it and compares it with $H$. Arithmetic coding, used in modern
compressors, reaches $H$ per symbol without forming blocks.
\end{remark}

\section{Typical sequences}

Theorem~\ref{thm:sourcecoding} measures entropy by codes of variable length. The asymptotic equipartition property gives
a more striking picture: among the $|A|^n$ possible sequences, almost all the probability is carried by about $2^{nH}$ of
them, which are all roughly equally likely.

Fix a finite alphabet $A$ and a law $p$ with $p(x)>0$ for all $x\in A$ (drop the other letters), let $X_1,X_2,\ldots$ be
independent with law $p$ and $H=H(X_1)$. For a sequence $x=(x_1,\dots,x_n)\in A^n$, write $p(x)=p(x_1)\cdots p(x_n)$, the probability
that $(X_1,\dots,X_n)=x$.

\begin{definition}
For $\varepsilon>0$, the \emph{typical set} $T_\varepsilon^{(n)}$ is the set of $x\in A^n$ with $2^{-n(H+\varepsilon)}\le p(x)\le2^{-n(H-\varepsilon)}$.
\end{definition}

\begin{theorem}[asymptotic equipartition property]\label{thm:aep}
Let $\varepsilon>0$.
\begin{enumerate}[label=(\alph*)]
\item The average surprise $-\frac1n\log_2p(X_1,\dots,X_n)$ tends to $H$ in probability; that is, the probability that
$(X_1,\dots,X_n)\in T_\varepsilon^{(n)}$ tends to $1$.
\item $|T_\varepsilon^{(n)}|\le2^{n(H+\varepsilon)}$ for all $n$, and $|T_\varepsilon^{(n)}|\ge(1-\varepsilon)2^{n(H-\varepsilon)}$ for all large $n$.
\item If $B_n\subseteq A^n$ has at most $2^{n(H-2\varepsilon)}$ elements, then $P\bigl((X_1,\dots,X_n)\in B_n\bigr)\to0$.
\end{enumerate}
\end{theorem}
\begin{proof}
(a) Let $Y_i=\log_2\frac1{p(X_i)}$. The $Y_i$ are independent (Proposition~\ref{prop:indepfun}), bounded, with mean $H$, and
$-\frac1n\log_2p(X_1,\dots,X_n)=\frac1n(Y_1+\dots+Y_n)$. By the weak law, it lies within $\varepsilon$ of $H$ with probability tending to $1$, and
this is exactly the event $(X_1,\dots,X_n)\in T_\varepsilon^{(n)}$.

(b) Every element of $T=T_\varepsilon^{(n)}$ has probability at least $2^{-n(H+\varepsilon)}$, so $1\ge P(T)\ge|T|2^{-n(H+\varepsilon)}$. Every element has
probability at most $2^{-n(H-\varepsilon)}$, so $P(T)\le|T|2^{-n(H-\varepsilon)}$, and $P(T)\ge1-\varepsilon$ for large $n$ by (a).

(c) Split $B_n$ into its typical and atypical elements:
$P(B_n)\le|B_n|\,2^{-n(H-\varepsilon)}+P\bigl((X_1,\dots,X_n)\notin T_\varepsilon^{(n)}\bigr)\le2^{-n\varepsilon}+o(1)$.
\end{proof}

\begin{figure}[ht]
\centering
\fig{ch04-aep}
\caption{The law of $-\frac1n\log_2p(X_1,\dots,X_n)$ for independent Bernoulli$(0.2)$ bits, for $n=20$, $100$ and $500$ (computed exactly:
the value only depends on the number of $1$s). It concentrates at $H=h(0.2)\approx0.722$ (dashed). The most likely sequence, all
zeros, has $-\frac1n\log_2p=0.32$ and is not typical: there is only one such sequence, while there are about $2^{0.722n}$ typical ones.}
\label{fig:aep}
\end{figure}

\begin{corollary}[fixed-length compression]\label{cor:fixedlength}
Let $R>0$ and $\mathbf X_n=(X_1,\dots,X_n)$. If $R>H$, there are maps $e_n:A^n\to\{0,1\}^{\lceil nR\rceil}$ and $d_n:\{0,1\}^{\lceil nR\rceil}\to A^n$
with $P\bigl(d_n(e_n(\mathbf X_n))\ne\mathbf X_n\bigr)\to0$. If $R<H$, every such pair has $P\bigl(d_n(e_n(\mathbf X_n))=\mathbf X_n\bigr)\to0$.
\end{corollary}
\begin{proof}
If $R>H$, let $\varepsilon=R-H$. By the theorem, $T_\varepsilon^{(n)}$ has at most $2^{nR}$ elements, so it can be mapped injectively into
$\{0,1\}^{\lceil nR\rceil}$; let $e_n$ be such a map on $T_\varepsilon^{(n)}$, arbitrary elsewhere, and $d_n$ invert it on the image. Decoding fails only
outside $T_\varepsilon^{(n)}$, with probability tending to $0$. If $R<H$, let $\varepsilon=\frac{H-R}3$. The set $B_n$ of sequences that are
decoded correctly satisfies $d_n(e_n(x))=x$, so $e_n$ is injective on $B_n$ and $|B_n|\le2^{\lceil nR\rceil}\le2^{nR+1}\le2^{n(H-2\varepsilon)}$ as soon as
$n\varepsilon\ge1$. By (c), $P(B_n)\to0$.
\end{proof}

So $H$ is a sharp threshold: compress below the entropy and you fail almost surely, above it and you succeed
almost surely. A text with $H$ bits of entropy per letter is, for the purpose of storage, a uniformly random
string of $2^{nH}$ possibilities, and the rest of the $|A|^n$ sequences may as well not exist.

\section{Exercises}

\begin{exercise}[$\star$]\label{ex:sharp}
(a) For $a>0$, find a random variable $X\ge0$, not constant, with $P(X\ge a)=\frac{\mathbb E[X]}a$. (b) For $a>0$ and $0<\sigma\le a$, find $X$ with mean
$0$, variance $\sigma^2$ and $P(|X|\ge a)=\frac{\sigma^2}{a^2}$.
\end{exercise}

\begin{exercise}[$\star$]\label{ex:lipschitz}
Let $f:[0,1]\to\R$ satisfy $|f(x)-f(y)|\le L|x-y|$. Show that $|B_nf(x)-f(x)|\le\frac L{2\sqrt n}$ for all $x$ and $n$.
\end{exercise}

\begin{exercise}[$\star\star$]\label{ex:pinsker}
Show that $I_p(a)\ge2(a-p)^2$ for $0<p<1$ and $0\le a\le1$, and deduce the upper bound of Theorem~\ref{thm:cramer} from Hoeffding's
inequality in the weaker form $e^{-2n(a-p)^2}$. When are the two exponents close?
\end{exercise}

\begin{exercise}[$\star\star$]\label{ex:entropybinom}
For $0<\alpha\le\frac12$, show that $\sum_{0\le k\le\alpha n}\binom nk\le2^{nh(\alpha)}$, and that $\binom nk\ge\frac1{n+1}2^{nh(k/n)}$ for $0\le k\le n$. So the number of
binary words of length $n$ with $\alpha n$ ones is $2^{nh(\alpha)+o(n)}$.
\end{exercise}

\begin{exercise}[$\star\star$]\label{ex:ballsbins}
Throw $n$ balls independently and uniformly into $n$ bins. (a) Show that a given bin receives at least $k$ balls with
probability at most $\binom nkn^{-k}\le\frac1{k!}\le(\frac ek)^k$. (b) Deduce that, for $n\ge3$, the fullest bin has fewer than
$k_n=\lceil3\ln n/\ln\ln n\rceil$ balls with probability at least $1-\frac1n$.
\end{exercise}

\begin{exercise}[$\star\star$]\label{ex:bloom}
A \emph{Bloom filter} stores a set of $n$ keys in an array of $m$ bits, initially $0$: each key sets to $1$ the bits at $k$
positions chosen by hash functions, modelled as independent and uniform. A key is reported present if its $k$ bits
are all $1$. (a) Show that a given bit is still $0$ after the $n$ insertions with probability $(1-\frac1m)^{kn}\le e^{-kn/m}$. (b) Assume that,
given the filter, a key not in the set has its $k$ positions independent and uniform, and that the fraction of zero bits is
close to its mean $e^{-kn/m}$. Its probability of being reported present (a \emph{false positive}) is then about
$(1-e^{-kn/m})^k$. Show that, for fixed $m/n$, this is smallest for $k=\frac mn\ln2$, where it equals $2^{-k}\approx0.6185^{m/n}$.
\end{exercise}

\begin{exercise}[$\star\star$]\label{ex:median}
\emph{The median trick.} A randomised algorithm returns, with probability at least $\frac34$, a number within $\varepsilon$ of an unknown
$\theta$. Run it $2k+1$ times independently and output the median of the answers. Show that the median is within $\varepsilon$ of $\theta$ with
probability at least $1-e^{-(2k+1)/8}$. How many runs give a failure probability below $10^{-6}$?
\end{exercise}

\begin{exercise}[$\star\star$]\label{ex:mcmillan}
A code is \emph{uniquely decodable} if any finite concatenation of codewords determines the sequence of symbols. (a) Show
that a prefix code is uniquely decodable, and that $\{0,01\}$ is uniquely decodable but not prefix. (b) Prove McMillan's
theorem: a uniquely decodable code for a finite alphabet satisfies Kraft's inequality, so Theorem~\ref{thm:sourcecoding}(a) holds for
it. (Expand $(\sum_x2^{-\ell(x)})^N$.)
\end{exercise}

\begin{exercise}[$\star\star$]\label{ex:subadd}
For random variables $X$, $Y$ with finitely many values, show that $H(X,Y)\le H(X)+H(Y)$, with equality iff $X$ and $Y$ are
independent. Interpret: describing a pair costs at most the cost of describing each.
\end{exercise}

\begin{exercise}[$\star\star\star$]\label{ex:longestrun}
Let $L_n$ be the length of the longest run of heads among the first $n$ tosses of a fair coin. Show that $\frac{L_n}{\log_2n}\to1$
almost surely. (For the lower bound, cut the first $n$ tosses into blocks of length $r=\lceil(1-\varepsilon)\log_2n\rceil$.)
\end{exercise}

\begin{exercise}[Journey]\label{pb:capset}
\emph{The cap set theorem.} Let $A\subset\mathbb F_3^n$ contain no three distinct points $x,y,z$ with $x+y+z=0$. We prove
$|A|\le3\cdot2.756^n$, following Ellenberg--Gijswijt and Tao.

For a finite set $A$ and a field $F$, a function $T:A^3\to F$ is a \emph{slice} if it has one of the forms $f(x)g(y,z)$,
$f(y)g(x,z)$, $f(z)g(x,y)$. The \emph{slice rank} of $T$ is the least number of slices whose sum is $T$.
\begin{enumerate}[label=(\alph*)]
\item Show that every subspace of $F^A$ of dimension $m$ contains a vector with at least $m$ non-zero coordinates.
\item Let $T(x,y,z)=\sum_{a\in A}c_a\,[x=y=z=a]$ with every $c_a\ne0$. Show that the slice rank of $T$ is $|A|$.
\item Show that for $x,y,z\in A$, $x+y+z=0$ if and only if $x=y=z$. Deduce that $T(x,y,z)=[x+y+z=0]$ has slice rank $|A|$.
\item Write $[x+y+z=0]=\prod_{i=1}^n\bigl(1-(x_i+y_i+z_i)^2\bigr)$ in $\mathbb F_3$, and deduce that the slice rank is at most $3M_n$,
where $M_n$ is the number of $\alpha\in\{0,1,2\}^n$ with $\alpha_1+\dots+\alpha_n\le\frac{2n}3$.
\item Show that $M_n\le\theta^n$ with $\theta=\min_{0<t\le1}\frac{1+t+t^2}{t^{2/3}}$, compute $\theta=2.7551\ldots$, and
conclude.
\end{enumerate}
\end{exercise}

\begin{exercise}[$\star\star\star$]\label{pb:kac}
\emph{Lacunary sums look Gaussian.} Let $t$ be uniform on $[0,2\pi]$ and, for $N\ge1$,
\[
S_N=\sum_{k=1}^N\cos(2^kt),\qquad T_N=\sum_{k=1}^N\cos(kt).
\]
\begin{enumerate}[label=(\alph*)]
\item Show that $\mathbb E[S_N]=\mathbb E[T_N]=0$ and $\mathbb E[S_N^2]=\mathbb E[T_N^2]=\frac N2$.
\item Show that for $N\ge2$,
\[
\mathbb E[S_N^4]=\frac34N^2+\frac98N-3,
\]
so that $\mathbb E[S_N^4]/\mathbb E[S_N^2]^2\to3$, the fourth moment of a standard Gaussian.
\item Show that $\mathbb E[T_N^4]\sim\frac{N^3}3$, so that $\mathbb E[T_N^4]/\mathbb E[T_N^2]^2\to\infty$.
\end{enumerate}
\end{exercise}

\begin{exercise}[$\star\star\star$]\label{pb:bonami}
\emph{The Bonami lemma.} Let $f:\{-1,1\}^n\to\R$ be a multilinear polynomial of degree at most $d$.
\begin{enumerate}[label=(\alph*)]
\item Show that $\mathbb E[f^4]\le9^d\,\mathbb E[f^2]^2$.
\item Deduce that if $f\ne0$, then $P\bigl(|f(x)|\ge\frac12\|f\|_2\bigr)\ge\frac9{16}\,9^{-d}$, where $\|f\|_2=\mathbb E[f^2]^{1/2}$.
\item Show that $9$ cannot be replaced by any constant $C<3$: for every $d$ and every $C<3$ there is $f$ of degree $d$ with
$\mathbb E[f^4]>C^d\,\mathbb E[f^2]^2$.
\end{enumerate}
\end{exercise}

\begin{exercise}[$\star\star$]\label{pb:kacring}
\emph{Kac's ring.} Label the sites of a circle by $\Z/n\Z$, and the edge from $j$ to $j+1$ by $j$. Each edge is marked independently with
probability $\mu\in\left]0,1\right[$; put $\varepsilon_j=-1$ if edge $j$ is marked and $\varepsilon_j=1$ otherwise, and let $M$ be the number of marked edges. Each site carries a
ball of colour $c_i=\pm1$ ($+1$ for black). At each tick every ball moves from its site $j$ to $j+1$, and changes colour if edge $j$ is
marked. Let $\Delta_t$ be the sum of the colours (the number of black balls minus the number of white ones) after $t$ ticks.
\begin{enumerate}[label=(\alph*)]
\item Show that the dynamics can be run backwards, that $\Delta_n=(-1)^M\Delta_0$, and that after $2n$ ticks every ball is back on its site with
its initial colour.
\item Show that the ball that starts at $i$ has colour $c_i\varepsilon_i\varepsilon_{i+1}\cdots\varepsilon_{i+t-1}$ after $t$ ticks, and that $\mathbb E[\Delta_t]=(1-2\mu)^t\Delta_0$ for $0\le t\le n$.
\item Show that $\operatorname{Var}(\Delta_t)\le n(2t-1)$ for $1\le t\le\frac{n+1}2$. Deduce that, for fixed $t$ and initial proportion $\beta$ of black balls, the
proportion of black balls at time $t$ tends in probability to $\frac12+(1-2\mu)^t(\beta-\frac12)$ as $n\to\infty$.
\item Starting with all balls black, the proportion of black balls is close to $\frac12$ after a few dozen ticks, yet it is $0$ or $1$ at
tick $n$, and $1$ at tick $2n$. Explain why there is no contradiction, and answer the questions of the think box.
\end{enumerate}
\end{exercise}

\begin{exercise}[$\star\star$]\label{pb:nbonacci}
\emph{Fibonacci with $n$ terms.} For $n\ge2$ let $P_n(x)=x^n-x^{n-1}-\dots-x-1$.
\begin{enumerate}[label=(\alph*)]
\item Show that $P_n$ has exactly one positive root $r_n$, that $1<r_n<2$, and that $r_n<r_{n+1}$. What is $r_2$?
\item Show that $r_n=2-2^{-n}-n\,2^{-2n-1}+O(n^22^{-3n})$. (Use $(x-1)P_n(x)=x^{n+1}-2x^n+1$.)
\item The number of coin-toss sequences of length $\ell$ without $n$ consecutive heads satisfies $a_\ell=a_{\ell-1}+\dots+a_{\ell-n}$. Explain why it grows like
$r_n^\ell$, and why $r_n$ is so close to $2$.
\end{enumerate}
\end{exercise}

\section{Hints}

\begin{hint}{ex:sharp}
(a) Two values. (b) Three values, symmetric around $0$.
\end{hint}
\begin{hint}{ex:lipschitz}
$|B_nf(x)-f(x)|\le L\,\mathbb E|\frac{S_n}n-x|$, and use Cauchy--Schwarz.
\end{hint}
\begin{hint}{ex:pinsker}
Fix $p$ and study $\varphi(a)=I_p(a)-2(a-p)^2$: compute $\varphi(p)$, $\varphi'(p)$ and $\varphi''$.
\end{hint}
\begin{hint}{ex:entropybinom}
Upper bound: $1\ge P(S\le\alpha n)$ for $S$ binomial$(n,\alpha)$, and $\alpha^k(1-\alpha)^{n-k}\ge\alpha^{\alpha n}(1-\alpha)^{(1-\alpha)n}$ for $k\le\alpha n$. Lower bound: look at the
proof of Theorem~\ref{thm:cramer}.
\end{hint}
\begin{hint}{ex:ballsbins}
(a) Union bound over the sets of $k$ balls. Compare $k!$ with $k^ke^{-k}$ using the exponential series. (b) Show that $(\frac ek)^k\le n^{-2}$
for $k=k_n$, and use a union bound over the bins.
\end{hint}
\begin{hint}{ex:bloom}
(b) Write $\rho=e^{-kn/m}$, so $k=-\frac mn\ln\rho$, and minimise $\ln\bigl((1-\rho)^k\bigr)=-\frac mn\ln\rho\ln(1-\rho)$ over $\rho\in]0,1[$.
\end{hint}
\begin{hint}{ex:median}
If the median is bad, at least $k+1$ runs are bad. Use Hoeffding.
\end{hint}
\begin{hint}{ex:mcmillan}
(b) Each term of the expansion corresponds to a word of $N$ symbols; group them by the length of their encoding.
\end{hint}
\begin{hint}{ex:subadd}
$H(X)+H(Y)-H(X,Y)$ is a relative entropy.
\end{hint}
\begin{hint}{ex:longestrun}
Upper bound: Proposition~\ref{prop:runs}. Lower bound: at least $\lfloor n/r\rfloor$ disjoint blocks, each all heads with probability $2^{-r}$; bound the
probability that none is, and use Borel--Cantelli (a).
\end{hint}

\begin{hint}{pb:capset}
(b) Contract $T$ against a function $h(z)$ that is orthogonal to the $f$'s of the third kind of slices and has large support.
(d) In each monomial one of the three variables has degree at most $\frac{2n}3$.
\end{hint}

\begin{hint}{pb:kac}
(a) Products of cosines are sums of cosines. (b) $\mathbb E[S_N^4]$ is $\frac1{16}$ of the number of $(k_1,\dots,k_4,\varepsilon_1,\dots,\varepsilon_4)$ with
$\sum\varepsilon_i2^{k_i}=0$. By uniqueness of binary expansions, such a relation either cancels in pairs or reads $2^{m+2}=2^{m+1}+2^m+2^m$.
(c) The same count for $\sum\varepsilon_ik_i=0$ is a volume.
\end{hint}

\begin{hint}{pb:bonami}
(a) Induction on $n$: $f=x_ng+h$ with $g$, $h$ not involving $x_n$, $\deg g\le d-1$. Expand the fourth power and use Cauchy--Schwarz on the
middle term. (b) Paley--Zygmund: Cauchy--Schwarz on $\mathbb E[f^21_{f^2\ge s/4}]$. (c) Products of normalised sums over disjoint blocks.
\end{hint}

\begin{hint}{pb:kacring}
(b) The edges $i,\dots,i+t-1$ are distinct when $t\le n$. (c) Two balls whose windows of edges are disjoint have independent colours. (d) For which
$t$ does (c) say something?
\end{hint}

\begin{hint}{pb:nbonacci}
(a) $P_{n+1}(x)=xP_n(x)-1$. (b) With $x=2-\varepsilon$, the equation becomes $\varepsilon=(2-\varepsilon)^{-n}$. Get $\varepsilon\le(2-\varepsilon_2)^{-n}$ with $\varepsilon_2=2-r_2$, then iterate.
(c) The characteristic polynomial of the recurrence is $P_n$.
\end{hint}

\section{Solutions}

\begin{solution}{ex:sharp}
\idea Markov's proof loses only where $X\ne a\mathbf 1_{X\ge a}$, so put all the mass at $0$ and $a$.
\writeup (a) Let $P(X=a)=t$ and $P(X=0)=1-t$ with $0<t<1$. Then $\mathbb E[X]=at$ and $P(X\ge a)=t=\frac{\mathbb EX}a$. (b) Apply (a) to $X^2$: let
$P(X=\pm a)=\frac{\sigma^2}{2a^2}$ each and $P(X=0)=1-\frac{\sigma^2}{a^2}$. Then $\mathbb EX=0$, $\operatorname{Var}X=\mathbb E[X^2]=\sigma^2$ and $P(|X|\ge a)=\frac{\sigma^2}{a^2}$.
\end{solution}

\begin{solution}{ex:lipschitz}
\idea The error is controlled by the average distance between the frequency and $x$, which is about one standard
deviation.
\writeup $|B_nf(x)-f(x)|=\bigl|\mathbb E[f(\frac{S_n}n)-f(x)]\bigr|\le L\,\mathbb E|\frac{S_n}n-x|$. By Cauchy--Schwarz (Proposition~\ref{prop:cs}),
$\mathbb E|\frac{S_n}n-x|\le\bigl(\mathbb E[(\frac{S_n}n-x)^2]\bigr)^{1/2}=\sqrt{\frac{x(1-x)}n}\le\frac1{2\sqrt n}$. For $f(x)=|x-\frac12|$ ($L=1$) at $x=\frac12$ the error is $\mathbb E|\frac{S_n}n-\frac12|$,
which is of order $\frac1{\sqrt n}$ by the central limit theorem of Chapter~6, so the rate cannot be improved in general.
\end{solution}

\begin{solution}{ex:pinsker}
\idea Both sides vanish to second order at $a=p$, and the second derivative of $I_p$ is at least $4$.
\writeup Let $\varphi(a)=I_p(a)-2(a-p)^2$ on $]0,1[$. Then $\varphi(p)=0$, $\varphi'(a)=\ln\frac ap-\ln\frac{1-a}{1-p}-4(a-p)$, so $\varphi'(p)=0$, and
$\varphi''(a)=\frac1a+\frac1{1-a}-4=\frac1{a(1-a)}-4\ge0$ since $a(1-a)\le\frac14$. So $\varphi$ is convex with a critical point at $p$, hence $\varphi\ge\varphi(p)=0$
(Volume~0, Proposition~4.1(b)); by continuity this extends to $a\in\{0,1\}$. Hence $e^{-nI_p(a)}\le e^{-2n(a-p)^2}$, and the latter is also Hoeffding's
bound with $t=n(a-p)$. The exponents are close when $\varphi''$ is close to $0$ between $p$ and $a$, that is, when $a$ and $p$ are both near $\frac12$:
for $p=\frac12$, $a=0.6$ they are $0.0201$ and $0.0200$, while for $p=0.1$, $a=0.2$ they are $0.044$ and $0.020$ (Figure~\ref{fig:chernoff}).
\end{solution}

\begin{solution}{ex:entropybinom}
\idea A binomial$(n,\alpha)$ law gives each word with $k$ ones the same weight, and the weights add up to at most $1$.
\writeup Let $S$ be binomial$(n,\alpha)$. For $k\le\alpha n$, since $\frac\alpha{1-\alpha}\le1$, $\alpha^k(1-\alpha)^{n-k}=(1-\alpha)^n(\frac\alpha{1-\alpha})^k\ge(1-\alpha)^n(\frac\alpha{1-\alpha})^{\alpha n}=2^{-nh(\alpha)}$. So
\[
\begin{aligned}1\ge P(S\le\alpha n)&=\sum_{k\le\alpha n}\binom nk\alpha^k(1-\alpha)^{n-k}\eqbr\ge2^{-nh(\alpha)}\sum_{k\le\alpha n}\binom nk.\end{aligned}
\]
For the lower bound, with $\beta=\frac kn$, the proof of Theorem~\ref{thm:cramer} showed that $\binom nk\beta^k(1-\beta)^{n-k}\ge\frac1{n+1}$ (the largest of
$n+1$ probabilities adding up to $1$), and $\beta^k(1-\beta)^{n-k}=2^{-nh(\beta)}$. (For $k\in\{0,n\}$, both sides are trivial.)
\end{solution}

\begin{solution}{ex:ballsbins}
\idea A bin is full only if some set of $k$ balls all land in it; Stirling-type bounds on $k!$ do the rest.
\writeup (a) If bin $1$ receives at least $k$ balls, some set of $k$ balls all land in it. There are $\binom nk$ such sets, each landing
entirely in bin $1$ with probability $n^{-k}$; the union bound gives $\binom nkn^{-k}\le\frac{n^k}{k!}n^{-k}=\frac1{k!}$. Finally $e^k=\sum_j\frac{k^j}{j!}\ge\frac{k^k}{k!}$, so
$\frac1{k!}\le(\frac ek)^k$.
(b) Let $L=\ln n$ and $k=k_n\ge\frac{3L}{\ln L}$. Then $\ln k\ge\ln3+\ln L-\ln\ln L$, and since $k(\ln k-1)$ increases with $k$ (for $k\ge1$),
\[
\begin{aligned}k(\ln k-1)&\ge\frac{3L}{\ln L}\bigl(\ln L-\ln\ln L+\ln3-1\bigr)\eqbr=3L\Bigl(1-\frac{\ln\ln L+1-\ln3}{\ln L}\Bigr).\end{aligned}
\]
With $u=\ln L>0$ (that is, $n\ge3$), the fraction is $\psi(u)=\frac{\ln u+1-\ln3}u$, and $\psi'(u)=\frac{\ln3-\ln u}{u^2}$, so its maximum is $\psi(3)=\frac13$. Hence
$k(\ln k-1)\ge2L$, i.e.\ $(\frac ek)^k\le n^{-2}$. By the union bound over the $n$ bins, some bin has at least $k_n$ balls with
probability at most $n\cdot n^{-2}=\frac1n$. (A matching lower bound of order $\frac{\ln n}{\ln\ln n}$ also holds, see Gonnet (1981); numerical evaluation shows the slow growth.)
\end{solution}

\begin{solution}{ex:bloom}
\idea Each of the $kn$ hash values misses a given bit with probability $1-\frac1m$; then optimise a function of one
variable after the right change of variable.
\writeup (a) The bit stays $0$ iff all $kn$ independent uniform positions avoid it: probability $(1-\frac1m)^{kn}\le e^{-kn/m}$ by $1-u\le e^{-u}$.
(b) If a fraction $\rho$ of the bits are $0$, each of the $k$ independent positions of a new key hits a $1$ with probability $1-\rho$, so
the false positive probability is $(1-\rho)^k$. With $\rho=e^{-kn/m}$, $k=-\frac mn\ln\rho$ and
\[
\ln\bigl((1-\rho)^k\bigr)=-\frac mn\ln\rho\,\ln(1-\rho).
\]
We must maximise $g(\rho)=\ln\rho\,\ln(1-\rho)>0$ on $]0,1[$. Its derivative is
\[
\begin{aligned}g'(\rho)&=\frac{\ln(1-\rho)}\rho-\frac{\ln\rho}{1-\rho}\eqbr=-\frac{\varphi(\rho)}{\rho(1-\rho)},\end{aligned}
\]
where $\varphi(\rho)=\rho\ln\rho-(1-\rho)\ln(1-\rho)$.
Now $\varphi(1-\rho)=-\varphi(\rho)$, $\varphi(\frac12)=0$, $\varphi(\rho)\to0$ as $\rho\to0$, and $\varphi'(\rho)=\ln\bigl(\rho(1-\rho)\bigr)+2$ increases on $]0,\frac12]$. So $\varphi'$ changes sign
at most once on $]0,\frac12]$, from negative to positive: $\varphi$ first decreases, then increases up to $\varphi(\frac12)=0$. Starting from the limit $0$
and ending at $0$, it is negative on $]0,\frac12[$, and by antisymmetry positive on $]\frac12,1[$. So $g'>0$ before $\frac12$ and $g'<0$ after:
the optimum is $\rho=\frac12$. Half the bits are $0$, $k=\frac mn\ln2$, and the
false positive rate is $2^{-k}=(2^{-\ln2})^{m/n}\approx0.6185^{m/n}$. With $10$ bits per key, $k\approx7$ and about $0.8\%$ false positives.
\end{solution}

\begin{solution}{ex:median}
\idea The median is wrong only if at least half the runs are wrong, which is a large deviation for a sum of
independent indicators of mean at most $\frac14$.
\writeup Let $Z_i$ be the indicator that run $i$ is bad (outside $[\theta-\varepsilon,\theta+\varepsilon]$), so $P(Z_i=1)\le\frac14$, and $Z=\sum Z_i$. If at most $k$
runs are bad, at least $k+1$ of the $2k+1$ answers lie in $[\theta-\varepsilon,\theta+\varepsilon]$, and so does the median (the $(k+1)$-th smallest
answer has at most $k$ answers below it, so it cannot lie below the interval, and symmetrically). So a bad median
forces $Z\ge k+1$, and $\mathbb EZ\le\frac{2k+1}4$, so $Z-\mathbb EZ\ge k+1-\frac{2k+1}4\ge\frac{2k+1}4$. Hoeffding with $n=2k+1$ and $t=\frac n4$ gives
$P(\text{bad median})\le e^{-2(n/4)^2/n}=e^{-n/8}$. For $10^{-6}$ we need $n\ge8\ln10^6\approx110.5$, so $111$ runs: the cost of a reliable
answer is only logarithmic in the confidence.
\end{solution}

\begin{solution}{ex:mcmillan}
\idea If a code is uniquely decodable, there are at most $2^L$ words of $N$ symbols whose encoding has length $L$; raising
Kraft's sum to a large power turns this count into the inequality.
\writeup (a) Decode a prefix code left to right: the first codeword is the unique codeword that is a prefix of the
string, remove it and repeat. The code $\{0,01\}$ is not prefix, but it is uniquely decodable: every $1$ is preceded by a $0$
that belongs to it, so reading right to left decodes. (b) Let $K=\sum_x2^{-\ell(x)}$ and $\ell_{\max}=\max\ell$. Then
\[
K^N=\sum_{x_1,\dots,x_N}2^{-(\ell(x_1)+\dots+\ell(x_N))}=\sum_{L=1}^{N\ell_{\max}}a_L2^{-L},
\]
where $a_L$ is the number of words $x_1\cdots x_N$ whose encoding has length $L$. These encodings are distinct binary strings of length
$L$ by unique decodability, so $a_L\le2^L$ and $K^N\le N\ell_{\max}$. If $K>1$, $K^N$ grows exponentially, a contradiction for large $N$. So
$K\le1$, and the proof of Theorem~\ref{thm:sourcecoding}(a) only used Kraft's inequality.
\end{solution}

\begin{solution}{ex:subadd}
\idea The defect of additivity compares the joint law with the product of its marginals.
\writeup Let $r(x,y)=P(X=x,Y=y)$, $p$ and $q$ the marginal laws. On the set where $r>0$ (where also $p,q>0$),
$\log_2\frac1{p(x)}+\log_2\frac1{q(y)}-\log_2\frac1{r(x,y)}=\log_2\frac{r(x,y)}{p(x)q(y)}$. Multiply by $r(x,y)$ and sum; since $\sum_yr(x,y)=p(x)$ and $\sum_xr(x,y)=q(y)$,
\[
H(X)+H(Y)-H(X,Y)=D(r\,\|\,p\otimes q)\ge0,
\]
where $p\otimes q(x,y)=p(x)q(y)$. By Gibbs, equality holds iff $r=p\otimes q$, i.e.\ iff $X$ and $Y$ are independent. This difference is the
\emph{mutual information} of $X$ and $Y$: the number of bits that one variable tells about the other.
\end{solution}

\begin{solution}{ex:longestrun}
\idea The longest run up to $n$ is at least the run at a single place and at most the longest of the runs starting
before $n$; Borel--Cantelli handles both sides, the lower one with disjoint blocks.
\writeup \emph{Upper bound.} $L_n\le\max_{m\le n}\ell_m$. By Proposition~\ref{prop:runs}, almost surely there is $m_0$ (random) with $\ell_m\le(1+\varepsilon)\log_2m$
for $m\ge m_0$. So for $n\ge m_0$, $L_n\le\max\bigl(\max_{m<m_0}\ell_m,(1+\varepsilon)\log_2n\bigr)$, and the first term is a finite random number (each $\ell_m$ is
finite almost surely, since $P(\ell_m\ge k)=2^{-k}\to0$). Hence $\limsup\frac{L_n}{\log_2n}\le1+\varepsilon$.
\emph{Lower bound.} Let $r=r_n=\lceil(1-\varepsilon)\log_2n\rceil$. The first $n$ tosses contain $\lfloor n/r\rfloor$ disjoint blocks of length $r$, and $L_n<r$ forces every
block to contain a tail. By independence of the blocks,
\[
\begin{aligned}P(L_n<r)&\le(1-2^{-r})^{\lfloor n/r\rfloor}\eqbr\le\exp\Bigl(-2^{-r}\Bigl(\frac nr-1\Bigr)\Bigr).\end{aligned}
\]
Now $2^{-r}\ge\frac12n^{-(1-\varepsilon)}$, so the exponent is at least $\frac{n^\varepsilon}{2r}-1$, with $r\le\log_2n+1$, and $\exp\bigl(-\frac{n^\varepsilon}{2(\log_2n+1)}\bigr)$ is summable (it is
eventually below $n^{-2}$). By Borel--Cantelli (a), almost surely $L_n\ge r_n\ge(1-\varepsilon)\log_2n$ for all large $n$. Intersecting over
$\varepsilon=\frac1k$ gives $\frac{L_n}{\log_2n}\to1$ almost surely.
\end{solution}

\begin{solution}{pb:capset}
\idea A diagonal tensor with $|A|$ non-zero entries has slice rank $|A|$, just as a diagonal matrix has rank $|A|$. A cap makes
$[x+y+z=0]$ diagonal on $A^3$. But the same function is a polynomial of degree $2n$, and in every monomial one variable has degree
at most $\frac{2n}3$: that writes it with few slices.

\writeup (a) Let $W$ have dimension $m$, and let $h\in W$ have the largest support $S$. If $|S|<m$, the subspace of $W$ of vectors
vanishing on $S$ is cut out by $|S|$ linear conditions, so it has dimension at least $m-|S|>0$ and contains some $w\ne0$. Since $w$
vanishes on $S$ and $h$ vanishes outside $S$, the support of $h+w$ is $S\cup\operatorname{supp}w\supsetneq S$, a contradiction.

(b) $T$ is the sum of the $|A|$ slices $c_a[x=a]\cdot[y=z=a]$, so its slice rank is at most $|A|$. Suppose
\[
T=\sum_{i=1}^{r_1}f_i(x)g_i(y,z)+\sum_{i=1}^{r_2}f'_i(y)g'_i(x,z)+\sum_{i=1}^{r_3}f''_i(z)g''_i(x,y).
\]
The functions $h:A\to F$ with $\sum_zf''_i(z)h(z)=0$ for all $i$ form a subspace of dimension at least $|A|-r_3$; by (a) one of
them has at least $|A|-r_3$ non-zero values. Multiply by $h(z)$ and sum over $z$. The left side becomes the diagonal matrix
$(x,y)\mapsto c_xh(x)[x=y]$, of rank at least $|A|-r_3$. The third sum disappears, and the first two become
$\sum_if_i(x)\tilde g_i(y)+\sum_i\tilde g'_i(x)f'_i(y)$, of rank at most $r_1+r_2$. So $r_1+r_2+r_3\ge|A|$.

(c) If $x+y+z=0$ and $x=y$, then $z=-2x=x$. If $x,y,z$ are pairwise distinct the cap condition forbids $x+y+z=0$. And if exactly
two coincide, say $x=y\ne z$, the first remark shows $x+y+z\ne0$. So on $A^3$, $[x+y+z=0]=[x=y=z]$, which is the $T$ of (b) with
$c_a=1$ (here $F=\mathbb F_3$).

(d) For $t\in\mathbb F_3$, $1-t^2$ is $1$ if $t=0$ and $0$ otherwise. Expanding the product gives a combination of monomials
$x^\alpha y^\beta z^\gamma$ with $\alpha,\beta,\gamma\in\{0,1,2\}^n$ and $|\alpha|+|\beta|+|\gamma|\le2n$, where
$|\alpha|=\sum\alpha_i$. In each monomial, one of $|\alpha|,|\beta|,|\gamma|$ is at most $\frac{2n}3$. Gather the monomials with
$|\alpha|\le\frac{2n}3$ by the value of $\alpha$: they form $\sum_\alpha x^\alpha G_\alpha(y,z)$, at most $M_n$ slices. Then those
remaining with $|\beta|\le\frac{2n}3$, then the rest, with $|\gamma|\le\frac{2n}3$. So the slice rank is at most $3M_n$, and by (c),
$|A|\le3M_n$.

(e) For $0<t\le1$, each $\alpha$ with $|\alpha|\le\frac{2n}3$ satisfies $t^{|\alpha|-2n/3}\ge1$, so
\[
M_n\le\sum_{\alpha\in\{0,1,2\}^n}t^{|\alpha|-2n/3}=\Bigl(\frac{1+t+t^2}{t^{2/3}}\Bigr)^n.
\]
The logarithmic derivative of $\frac{1+t+t^2}{t^{2/3}}$ vanishes when $\frac{1+2t}{1+t+t^2}=\frac2{3t}$, that is $4t^2+t-2=0$,
$t=\frac{\sqrt{33}-1}8=0.5931$, where the value is $\theta=2.7551$. So $|A|\le3\theta^n<3\cdot2.756^n$.

\tools Volume~0, Chapter~7 (rank, dimension); Volume~3, Chapter~4 (the exponential moment trick of Chernoff in (e)).
\leads The bound $\theta$ is optimal for this method, and the true growth rate is unknown: the best constructions give caps of size
about $2.22^n$ (Edel 2004; in 2023 a computer search guided by a language model, FunSearch, found a cap of $512$ points in $\mathbb F_3^8$).
The known maximum caps for $n=1,\dots,6$ are $2,4,9,20,45,112$, against $3M_n=3,9,30,45,153,504$. The slice rank method has since
proved the sunflower-free conjecture of Erdős and Szemerédi and bounded many other configurations.
\end{solution}

\begin{solution}{pb:kac}
\idea Moments of trigonometric sums count solutions of linear equations in the frequencies. With frequencies $2^k$, binary
expansions are unique and almost every solution cancels in pairs, exactly as for independent variables; that is why lacunary sums
behave like sums of independent terms. With frequencies $1,2,\dots,N$ there are far too many additive coincidences.

\writeup (a) Since $\cos A\cos B=\frac12(\cos(A-B)+\cos(A+B))$ and $\mathbb E[\cos mt]$ is $1$ for $m=0$ and $0$ for every other integer $m$, we get
$\mathbb E[\cos(at)\cos(bt)]=\frac12\cdot1_{a=b}$ for positive integers $a$, $b$. So both means vanish and both second moments are $\frac N2$.

(b) Expanding $\cos A_i=\frac12(e^{iA_i}+e^{-iA_i})$,
\[
\mathbb E\prod_{i=1}^4\cos(n_it)=\frac1{16}\#\Bigl\{\varepsilon\in\{\pm1\}^4:\sum_i\varepsilon_in_i=0\Bigr\},
\]
so $\mathbb E[S_N^4]=\frac M{16}$, where $M$ counts the $(k,\varepsilon)\in\{1,\dots,N\}^4\times\{\pm1\}^4$ with $\sum\varepsilon_i2^{k_i}=0$. In such a relation the terms with sign
$+$ and those with sign $-$ have equal sums; neither group is empty.

\emph{Two against two}, $2^a+2^b=2^c+2^d$. If $a\ne b$ the left side has exactly two binary digits $1$, at $a$ and $b$; the right side
then cannot be $2^{c+1}$, so $c\ne d$ and $\{c,d\}=\{a,b\}$. If $a=b$ the left side is $2^{a+1}$, which forces $c=d=a$. Either way the terms
cancel in pairs.

\emph{One against three}, $2^a=2^b+2^c+2^d$ with $b\le c\le d$. If $b,c,d$ are distinct the right side has three digits $1$; if they are
equal it is $3\cdot2^b$; if $b<c=d$ it is $2^b+2^{c+1}$ with two digits. So $b=c<d$, and $2^{b+1}+2^d$ is a power of $2$ only if $d=b+1$: the
relation is $2^{m+2}=2^{m+1}+2^m+2^m$, with $1\le m\le N-2$.

Counting. Pairings: choose one of the $3$ ways of splitting the four positions into two pairs, an exponent for each pair ($N^2$
ways) and the sign of the first element of each pair ($4$ ways): $12N^2$. A tuple is counted more than once only if it fits two
splittings, which happens exactly when the four exponents are equal (two signs $+$, two signs $-$: $6N$ tuples, each fitting $2$
splittings); so there are $12N^2-6N$ pairing tuples. The other relations: $N-2$ values of $m$, $\frac{4!}{2!}=12$ arrangements of the exponents
$m+2,m+1,m,m$, and $2$ sign choices: $24(N-2)$ tuples, none of which cancels in pairs. So
\[
M=12N^2-6N+24(N-2)=12N^2+18N-48,\qquad\mathbb E[S_N^4]=\frac34N^2+\frac98N-3,
\]
and $\mathbb E[S_N^4]/\mathbb E[S_N^2]^2=3+\frac9{2N}-\frac{12}{N^2}\to3$.

(c) Now count the solutions of $\sum\varepsilon_ik_i=0$ with $1\le k_i\le N$. For each of the $6$ sign patterns with two $+$ and two $-$, the number of
solutions of $k_1+k_2=k_3+k_4$ is $\sum_{s=2}^{2N}r(s)^2$ with $r(s)=\min(s-1,2N+1-s)$, that is $2\sum_{j=1}^{N-1}j^2+N^2=\frac{2N^3+N}3$. For each of the $8$
patterns with three signs alike, the number of solutions of $k_1+k_2+k_3=k_4$ is $\sum_{s\le N}\binom{s-1}2=\binom N3$. So
\[
\mathbb E[T_N^4]=\frac1{16}\Bigl(2(2N^3+N)+8\binom N3\Bigr)\sim\frac1{16}\Bigl(4+\frac43\Bigr)N^3=\frac{N^3}3,
\]
and $\mathbb E[T_N^4]/\mathbb E[T_N^2]^2\sim\frac{4N}3\to\infty$.

\tools{Volume~0, Chapter~2 (Euler's formula); Volume~1, Chapter~1 (binary expansions); Volume~3, Chapter~4 (moments).}
\leads The same counting shows that every moment of $S_N/\sqrt{N/2}$ tends to the corresponding Gaussian moment, and the method of
moments then gives Kac's theorem (1946): $S_N/\sqrt{N/2}$ converges in law to $N(0,1)$, although the terms $\cos(2^kt)$ are not
independent. Salem and Zygmund (1947) extended it to $\sum\cos(n_kt)$ whenever $n_{k+1}/n_k\ge q>1$. Lacunarity alone is not enough for
other functions: Erdős and Fortet observed that for $f(t)=\cos t+\cos2t$ and $n_k=2^k-1$, the limit law of $\sum_kf(n_kt)/\sqrt N$ is a mixture of
Gaussians, because of the relations $2n_k+1=n_{k+1}$.
\end{solution}

\begin{solution}{pb:bonami}
\idea Peel off one variable: $f=x_ng+h$. The fourth moment splits into three terms, and the cross term is controlled by Cauchy--Schwarz
and induction; the constants are exactly what makes $9^d$ reproduce itself.

\writeup (a) Induction on $n$; for $n=0$, $f$ is a constant $c$ and $\mathbb E[f^4]=c^4=\mathbb E[f^2]^2$. Write $f=x_ng+h$ with $g$, $h$ multilinear in $x_1,\dots,x_{n-1}$,
$\deg g\le d-1$ and $\deg h\le d$ (if $d=0$, $g=0$). As $x_n$ is independent of $g$, $h$ with $\mathbb E[x_n]=\mathbb E[x_n^3]=0$ and $x_n^2=1$,
\[
\mathbb E[f^2]=\mathbb E[g^2]+\mathbb E[h^2],\qquad\mathbb E[f^4]=\mathbb E[h^4]+6\,\mathbb E[g^2h^2]+\mathbb E[g^4].
\]
By induction $\mathbb E[h^4]\le9^d\mathbb E[h^2]^2$ and $\mathbb E[g^4]\le9^{d-1}\mathbb E[g^2]^2$, and by Cauchy--Schwarz $\mathbb E[g^2h^2]\le\sqrt{\mathbb E[g^4]\mathbb E[h^4]}\le3^{2d-1}\mathbb E[g^2]\mathbb E[h^2]$.
Since $6\cdot3^{2d-1}=2\cdot9^d$,
\[
\mathbb E[f^4]\le9^d\bigl(\mathbb E[h^2]^2+2\mathbb E[g^2]\mathbb E[h^2]+\mathbb E[g^2]^2\bigr)=9^d\,\mathbb E[f^2]^2.
\]

(b) Let $s=\mathbb E[f^2]>0$ and $q=P(f^2\ge\frac s4)$. Then $s\le\frac s4+\mathbb E[f^21_{f^2\ge s/4}]\le\frac s4+\sqrt{\mathbb E[f^4]}\sqrt q$, so
$q\ge\frac{9s^2}{16\,\mathbb E[f^4]}\ge\frac9{16}9^{-d}$ by (a).

(c) Split $md$ variables into $d$ blocks $B_1,\dots,B_d$ of size $m$, let $L_j=m^{-1/2}\sum_{i\in B_j}x_i$, and $f=L_1\cdots L_d$, multilinear of degree $d$. The
$L_j$ are independent, $\mathbb E[L_j^2]=1$ and $\mathbb E[L_j^4]=\frac1{m^2}\bigl(m+3m(m-1)\bigr)=3-\frac2m$. So $\mathbb E[f^2]=1$ and $\mathbb E[f^4]=(3-\frac2m)^d>C^d$ once $\frac2m<3-C$.

\tools{Volume~3, Chapter~4 (Cauchy--Schwarz for expectations, moments of sums of signs).}
\leads Bonami proved this in 1970, in a study of Fourier coefficients on groups; Beckner and Gross (1975) found the sharp form,
\emph{hypercontractivity}: $\|T_\rho f\|_4\le\|f\|_2$ for $\rho\le\frac1{\sqrt3}$, where $T_\rho$ multiplies $\hat f(S)$ by $\rho^{|S|}$. Applied to the derivatives
$D_if$, it is the key to the theorem of Kahn, Kalai and Linial quoted in Problem~\ref{pb:influence}, to Friedgut's theorem that functions
of small total influence depend essentially on few variables, and to much of the analysis of Boolean functions in complexity
theory. Part~(c) only shows that the base is at least $3$, but $9$ is in fact the right base: by the central limit theorem, polynomials on
the cube approximate polynomials in Gaussian variables, and for the Hermite polynomials $H_d$ the $d$-th root of
$\mathbb E[H_d^4]/\mathbb E[H_d^2]^2$ tends to $9$ (slowly: one finds numerically $7.8$ at $d=30$).
\end{solution}

\begin{solution}{pb:kacring}
\idea Each ball's colour is a product of independent signs along the window of edges it has crossed. Products over overlapping windows
are correlated, over disjoint windows independent, and there are only $O(t)$ overlapping windows per ball: the law of large numbers wins
for every fixed $t$, while the recurrence at time $2n$ moves away as $n$ grows.

\writeup (a) Backwards, every ball moves from $j+1$ to $j$ and changes colour if edge $j$ is marked; this undoes a tick. After $n$ ticks every ball has
crossed each edge once, so it has changed colour $M$ times and $\Delta_n=(-1)^M\Delta_0$; after $2n$ ticks it is on its site again, having changed colour $2M$ times.

(b) The ball starting at $i$ crosses edges $i,i+1,\dots,i+t-1$ during the first $t$ ticks, hence its colour. For $t\le n$ these edges are distinct, so
$\mathbb E[\varepsilon_i\cdots\varepsilon_{i+t-1}]=(1-2\mu)^t$, and $\mathbb E[\Delta_t]=(1-2\mu)^t\sum_ic_i$.

(c) Let $W_i=\{i,\dots,i+t-1\}$ and $\Pi_i=\prod_{j\in W_i}\varepsilon_j$. If $W_i\cap W_{i'}=\emptyset$, $\Pi_i$ and $\Pi_{i'}$ are independent and uncorrelated. For $t\le\frac{n+1}2$, $W_{i'}$ meets $W_i$ only when
the cyclic distance from $i$ to $i'$ is less than $t$, which happens for $2t-1$ values of $i'$. As $|\Pi_i|=1$, every covariance is at most $1$ in absolute value. So
$\operatorname{Var}(\Delta_t)=\sum_{i,i'}c_ic_{i'}\operatorname{Cov}(\Pi_i,\Pi_{i'})\le n(2t-1)$. By Chebyshev, $\frac{\Delta_t}n-(1-2\mu)^t\frac{\Delta_0}n\to0$ in probability; with $\frac{\Delta_0}n=2\beta-1$, the proportion of black
balls $\frac12(1+\frac{\Delta_t}n)$ tends to $\frac12+(1-2\mu)^t(\beta-\frac12)$.

(d) Part (c) is a statement for each fixed $t$ as $n\to\infty$ (and remains useful for $t$ small compared with $n$), while the returns of (a) happen at
times $n$ and $2n$, which run away to infinity. The variance bound $n(2t-1)$ is useless for $t$ of order $n$; at $t=n$ all windows are the whole
circle, all products equal $(-1)^M$, and the fluctuations are as large as possible. This is Boltzmann's answer to Zermelo: the
recurrence time of a gas is astronomically long. Loschmidt's reversal is answered by (c) as well: running the ring backwards from a
typical state at time $t$ is the dynamics of the ring with mirrored marks, but from a very special state, the only one that leads
back to all black. In the think box: with $3$ marked edges out of $6$, after $6$ ticks every ball has changed colour $3$ times, so all balls are
white; after $12$ ticks they are all black again. With $1000$ sites and $\mu=\frac1{10}$, after $10$ ticks the proportion of black balls is close to
$\frac12(1+0.8^{10})\approx0.554$, and after $1000$ ticks it is exactly $0$ or $1$, according to the parity of $M$.

\tools{Volume~3, Chapter~4 (variance, Chebyshev's inequality); Volume~2, Chapter~1 (the rotation group of the ring).}
\leads Gottwald and Oliver's article \emph{Boltzmann's dilemma} (2009) works out the Kac ring in detail and uses it as an introduction to
statistical mechanics. Kac's own account is in his \emph{Probability and Related Topics in Physical Sciences} (1959).
\end{solution}

\begin{solution}{pb:nbonacci}
\idea Multiplying by $x-1$ makes the equation $x^n(2-x)=1$: the root is $2$ minus something exponentially small, and the equation for that small thing can be iterated.

\writeup (a) $P_n(x)=x^n\bigl(1-\sum_{j=1}^nx^{-j}\bigr)$, and the bracket increases strictly on $\left]0,\infty\right[$ from $-\infty$ to $1$: exactly one positive root $r_n$. $P_n(1)=1-n<0$ and $P_n(2)=2^n-(2^n-1)=1>0$, so
$1<r_n<2$. Since $P_{n+1}(x)=xP_n(x)-1$, $P_{n+1}(r_n)=-1<0$, so $r_{n+1}>r_n$. And $r_2=\frac{1+\sqrt5}2$.

(b) Since $(x-1)P_n(x)=x^{n+1}-2x^n+1$, $r=r_n$ satisfies $r^n(2-r)=1$. Let $\varepsilon=2-r\in\left]0,2-r_2\right]$. Then $\varepsilon=(2-\varepsilon)^{-n}\le r_2^{-n}$, so $n\varepsilon\to0$. Now
$(2-\varepsilon)^{-n}=2^{-n}\exp\bigl(-n\ln(1-\frac\varepsilon2)\bigr)=2^{-n}\bigl(1+\frac{n\varepsilon}2+O(n^2\varepsilon^2)\bigr)$. This first gives $\varepsilon=2^{-n}(1+O(n\varepsilon))=O(2^{-n})$, then
$\varepsilon=2^{-n}+n2^{-n-1}\varepsilon+O(n^22^{-3n})=2^{-n}+n2^{-2n-1}+O(n^22^{-3n})$.

(c) A sequence without $n$ consecutive heads ends with a tail preceded by $0,\dots,n-1$ heads, or is shorter than $n$; this gives the recurrence, whose characteristic polynomial is
$P_n$. With $r=r_n$, $\sum_{i=1}^nr^{-i}=1$. If $a_j\ge r^{j-n}$ for $n$ consecutive $j$, then $a_\ell\ge\sum_{i=1}^nr^{\ell-i-n}=r^{\ell-n}$; the initial values $a_j=2^j$ ($j<n$) satisfy this, so by
induction $a_\ell\ge r^{\ell-n}$ for all $\ell$. In the same way $a_\ell\le c\,r^\ell$ with $c=\max_{j<n}2^jr^{-j}$. So $a_\ell^{1/\ell}\to r_n$. And $r_n$ is close to $2$ because forbidding a run of $n$
heads costs little: among the $2^\ell$ sequences, a fraction about $e^{-\ell2^{-n-1}}$ avoids such runs, consistent with $(\frac{r_n}2)^\ell\approx(1-2^{-n-1})^\ell$.

\tools{Volume~0, Chapter~5 (sequences, recurrences); Volume~3, Chapter~4 (runs of heads).}
\leads For $n=3$, $r_3=1.8393$ is the tribonacci constant. The estimate $(1-2^{-n-1})^\ell$ is the heart of the Erd\H os--R\'enyi law of the longest run (Volume~3): runs of length $n$ appear once
$\ell2^{-n}$ is large, that is $n\approx\log_2\ell$.
\end{solution}

\section*{Sources and further reading}

\begin{itemize}
\item J.~Bernoulli, \emph{Ars Conjectandi}, Basel, 1713; English translation by E.~D.~Sylla, \emph{The Art of
Conjecturing}, Johns Hopkins, 2006. Part~IV contains the first law of large numbers.
\item S.~Bernstein, \emph{D\'emonstration du th\'eor\`eme de Weierstrass fond\'ee sur le calcul des probabilit\'es},
Comm.\ Soc.\ Math.\ Kharkov 13 (1912), 1--2.
\item H.~Chernoff, \emph{A measure of asymptotic efficiency for tests of a hypothesis based on the sum of
observations}, Annals of Mathematical Statistics 23 (1952), 493--507; W.~Hoeffding, \emph{Probability inequalities for
sums of bounded random variables}, Journal of the American Statistical Association 58 (1963), 13--30.
\item C.~E.~Shannon, \emph{A mathematical theory of communication}, Bell System Technical Journal 27 (1948), 379--423
and 623--656. The birth of information theory, still very readable.
\item T.~M.~Cover and J.~A.~Thomas, \emph{Elements of Information Theory}, 2nd ed., Wiley, 2006, Chapters~2, 3 and 5:
entropy, the AEP, Kraft, McMillan and Huffman.
\item M.~Mitzenmacher and E.~Upfal, \emph{Probability and Computing}, 2nd ed., Cambridge University Press, 2017,
Chapters~3--5: Chernoff bounds, balls into bins and Bloom filters.
\item P.~Erd\H{o}s and A.~R\'enyi, \emph{On a new law of large numbers}, Journal d'Analyse Math\'ematique 23 (1970),
103--111. Longest runs and their generalisations; M.~F.~Schilling, \emph{The longest run of heads}, College
Mathematics Journal 21 (1990), 196--207, is a delightful account.
\item G.~H.~Gonnet, \emph{Expected length of the longest probe sequence in hash code searching}, Journal of the ACM
28 (1981), 289--304. The maximum load of balls in bins, to first order.
\item B.~H.~Bloom, \emph{Space/time trade-offs in hash coding with allowable errors}, Communications of the ACM 13
(1970), 422--426.
\end{itemize}
